\section{Methodology}
\textbf{Geometric Anchor Pretraining (GAP)} is a pretraining strategy designed to extract object-centric geometric priors for visuomotor imitation learning directly from dense VFM's embeddings. GAP addresses the spatial overfitting commonly observed in data-scarce regimes ($N \le 50$) by explicitly supervising a coordinate-based adapter via a masked proxy task. %
We begin by formalizing the imitation learning setting and defining the spatial adapter architecture. We then introduce our method (Figure \ref{fig:method_overview}), focusing on the proxy task and losses, followed by a description of the downstream policy adaptation procedure.

\subsection{Imitation Learning with Diffusion Policies}
\label{subsec:background}
We focus on visuomotor policy learning in the Imitation Learning setting, where an agent learns control behaviors from expert demonstrations. The agent is provided with a dataset $\mathcal{D} = \{ \tau_i \}_{i=1}^N$, where each trajectory $\tau_i = \{ (o_t, a_t) \}_{t=0}^T$ consists of visual observations $o_t$ and corresponding actions $a_t$. The objective is to infer a policy $\pi_\theta(a_t | o_t)$ that reproduces the expert behavior.

In this work, we adopt \textbf{Diffusion Policies} \cite{chi2023diffusion}, which parameterize the action distribution $\pi_\theta(a_t | o_t)$ as a conditional denoising diffusion probabilistic model. The training objective minimizes the noise prediction error:
\begin{equation}
    \mathcal{L}_{diff} = \mathbb{E}_{\epsilon \sim \mathcal{N}, t \sim \mathcal{U}, \tau \sim \mathcal{D}} \left[ \| \epsilon - \epsilon_\theta(a_t^{(k)}, k, E(o_t)) \|_2^2 \right]
\end{equation}
where $E(o_t)$ is the visual embedding used to condition the denoising network $\epsilon_\theta$ at diffusion step $k$. While a wide range of Imitation Learning algorithms have been introduced, with no loss of generality, in this work, we experiment with \cite{chi2023diffusion} as the policy learning method for all the experiments to isolate the impact of the conditioning representation. GAP focuses entirely on the pre-training of the adapter between the frozen vision encoder and the policy $\pi_\theta$ to provide a robust, geometry-aware conditioning signal.

\subsection{The Spatial Adapter}
\label{subsec:arch}
To extract robust semantic features, we use a frozen pretrained backbone $f_\phi$ (e.g., ResNet-50, ViT-S, or ViT-B). To map visual features into precise spatial embeddings, $f_\phi$ is followed by a lightweight adapter module, which we denote as $f_A$ (see \autoref{fig:method_overview}). Specifically, $f_A$ first applies a shallow convolutional network (a $3\times3$ followed by a $1\times1$ convolution) to project the high-dimensional backbone features into $K$ spatial activation maps, denoted as $\Phi_t \in \mathbb{R}^{K \times h \times w}$. Finally, $f_A$ applies a Spatial Softmax (SS) \cite{finn2016deep} module to convert these maps into $K$ 2D spatial coordinates, which we define as our candidate \emph{keypoints} $P_t = \{p_{k,t}\}_{k=1}^K$:

\begin{equation}
    p_{k,t} = \sum_{x=1}^{w} \sum_{y=1}^{h} \begin{bmatrix} x \\ y \end{bmatrix} \frac{\exp(\Phi_{t, k, x, y})}{\sum_{x',y'} \exp(\Phi_{t, k, x', y'})}
\end{equation}
where $\Phi_{t, k, x, y}$ is the activation of the $k$-th feature channel at spatial location $(x,y)$. 

This mapping converts the visual learning paradigm into a state-based one, compressing dense feature maps into $K$ sparse 2D coordinates. When training policies for 200+ demonstrations, we observe that the learned keypoints naturally stick to specific objects and become reliable ``semantic trackers''. However, under extremely low-data regime, without explicit supervision, these keypoints tend to latch onto spurious visual cues rather than geometrically meaningful locations. %
This occurs because the downstream action regression loss ($\mathcal{L}_{diff}$) provides only weak, indirect spatial supervision. Forced to minimize training error with minimal data, the network takes a shortcut, anchoring to high-contrast static distractors (e.g., table textures) instead of complex object geometry.

\subsection{Geometric Anchor Pretraining (GAP)}
\label{subsec:gap_method}
To address this, GAP pretrains the spatial bottleneck $f_A$ on a single, cheap simulated proxy task. This aims to decouple geometric feature learning from action-mapping.

\textbf{Proxy Task.} 
The objective of our training is to align geometric keypoints with task-relevant objects in the scene, without any task-specific knowledge. Therefore, in principle we can use for pretraining any simple manipulation task or contact-rich motor babbling. In this paper, with no loss of generality, we experiment with the \texttt{LiftCube} task from Robomimic \cite{mandlekar2021matters}---
the simplest task available in the benchmark---in which a Franka Emika Panda robot is tasked to reach, grasp, and lift a randomly positioned cube on a plain \emph{white-background} table (see \autoref{fig:qualitative_small}-a). To evaluate proxy task invariance, we also experiment with a different pre-training task (PlaceSphere from ManiSkill).
Trajectories are generated automatically via a scripted controller, requiring no human teleoperation. We use 100 demonstrations of this proxy task, leveraging ground-truth object segmentation masks $\mathcal{M}_t$ provided by the simulator to supervise our loss. Crucially, no expert action labels $a_t$ are needed at any point during this phase. %

\begin{figure}[ht!]
\centering
\includegraphics[width=0.85\columnwidth]{pics/qualitative_hq.png}
\caption{\textbf{Qualitative Keypoint Transfer.} When pre-trained on a proxy task and then transferred to a new task and simulator, GAP allows for keeping a favorable geometric grounding even in zero-shot. \textbf{(a)} shows the keypoints placement on the Robomimic \texttt{LiftCube} task after pretraining. \textbf{(b)} and \textbf{(c)} show keypoints positioning when using the pre-trained visual encoder in zero-shot on a different task (of the same simulator) and on a different task and simulator (Maniskill), respectively. We show a third-person view on the left and a first-person view on the right.}
\label{fig:qualitative_small}
\vspace{-3pt}
\end{figure}

\textbf{GAP Spatial Objectives.}
We supervise the adapter $f_A$ using a multi-objective spatial loss that enforces object-centric, spatially distributed, and non-redundant keypoints over the robot gripper and interacting objects. This is achieved through three components, depicted in \autoref{fig:method_overview} and detailed below.


\subsubsection{Centroid Alignment ($\mathcal{L}_{center}$)}
To ensure keypoints ground themselves on the target object rather than background distractors, we minimize the distance between the predicted keypoint centroid $\bar{p}_t$ and the ground-truth mask centroid $c_t$:
\begin{equation}
    \mathcal{L}_{center} = \| \bar{p}_t - c_t \|_2^2 \quad \text{where} \quad \bar{p}_t = \frac{1}{K} \sum_{k=1}^K p_{k,t}
\end{equation}
and $c_t$ is computed via the spatial moments of the binary mask $\mathcal{M}_t$.

\subsubsection{Geometric Spread ($\mathcal{L}_{spread}$)}
To prevent the degenerate solution where all keypoints collapse precisely in the centroid (which drops orientation information), we enforce the spatial variance of the keypoints $\sigma_p$ to match the normalized object scale $\sigma_{target}$:
\begin{equation}
    \mathcal{L}_{spread} = \| \sigma_p - \sigma_{target} \|_2^2 \quad \text{where} \quad \sigma_p = \frac{1}{K} \sum_{k=1}^K \| p_{k,t} - \bar{p}_t \|_2
\end{equation}
The target scale is derived from the mask area $A_t = \sum \mathcal{M}_t$, approximated as $\sigma_{target} = 0.8 \times \sqrt{A_t / \pi}$, representing a proportional bounding radius.

\subsubsection{Keypoint Diversity ($\mathcal{L}_{div}$)}
Lastly, to maximize the structural information captured by the bottleneck, we penalize redundancy by enforcing a minimum separation margin $\delta_{min}$ between any pair of keypoints:
\begin{equation}
    \mathcal{L}_{div} = \frac{1}{K} \sum_{k=1}^K \left[ \max\left(0, \delta_{min} - \min_{j \neq k} \| p_{k,t} - p_{j,t} \|_2 \right) \right]^2
\end{equation}
This term encourages the network to discover the object's distinct geometric extremities, thereby generating highly informative \emph{Geometric Anchors}.

The final pre-training objective combines these terms to create a "push-pull" dynamic: keypoints are pulled onto the object ($\mathcal{L}_{center}$), but pushed outward to span its geometry ($\mathcal{L}_{spread}$) and away from one another ($\mathcal{L}_{div}$).
\begin{equation}
    \mathcal{L}_{GAP} = \lambda_c \mathcal{L}_{center} + \lambda_s \mathcal{L}_{spread} + \lambda_d \mathcal{L}_{div}
\end{equation}
In our implementation, we prioritize spatial coverage over strict centering by setting $\lambda_c = 0.3$, $\lambda_s = 0.5$, and $\lambda_d = 2.0$, with a diversity margin $\delta_{min} = 0.15$ (normalized image coordinates). We report an extensive ablation on the different objective components in Sec. \ref{ablation_losses}.

\indent\textbf{Object-Centric Keypoint Allocation.} While end-to-end policies require massive datasets to naturally develop entity-centric keypoints, GAP explicitly enforces this optimal behavior in low-data regimes. Given a pre-training scene with $M$ semantic entities, we partition the $K$ available keypoints into $M$ disjoint subsets: $P_t = \bigcup_{m=1}^{M} P_{t,m}$. The spatial regularization objectives ($\mathcal{L}_{GAP}$) are then applied independently to each subset using its corresponding mask. 
This $M$-way partition bridges the abstract semantic expressivity of Vision Foundation Models (VFMs) with the strict, object-centric geometric priors required for physical manipulation. Consequently, downstream fine-tuning does not need to learn object separation from scratch. When transitioning from a proxy task to novel, multi-object environments (e.g., \texttt{StackCube}, \texttt{SquareNut}), these pre-trained subsets deploy as independent semantic trackers. Having already internalized priors of centroid alignment, spatial spread, and extremity-seeking diversity, they rapidly re-anchor to novel geometries with minimal adaptation. \autoref{fig:qualitative_small} qualitatively demonstrates this robust, entity-separated initialization prior to any policy training.




\subsection{Downstream Policy Adaptation}
\label{subsec:downstream_adapt}
Following GAP, the regularized adapter and frozen encoder $E$ are transferred to the downstream tasks. The $M$ pre-trained keypoint subsets drastically reduce the policy's learning burden: subsets tracking persistent elements (e.g., the manipulator) require light adaptation, allowing the few-shot demonstrations to focus entirely on grounding the remaining keypoints to novel target objects. This preserves learned spatial priors and enables highly sample-efficient convergence.
For fair comparison, \textbf{all evaluated baselines}---including end-to-end models and VFMs (R3M, DINOv2, VC-1)---employ an identical architecture. Models denoted with ``SS'' utilize our convolutional pooler and Spatial Softmax. During downstream fine-tuning, privileged segmentation masks $\mathcal{M}_t$ are strictly discarded. The vision backbone remains frozen, and we fine-tune only the lightweight adapter $f_A$ alongside the diffusion head via the action-prediction objective $\mathcal{L}_{diff}$.
